% =============================================================
% Chapitre 1 — Contexte du projet, étude de l'existant et problématique (~9 pages)
% Source : ../../CH1_Cadre_general_et_etude_de_l_existant.md
% Enrichissement §1.1 et §1.2 : DOSSIER_TECHNIQUE_MENAL.md (§18.1 à §18.6)
% =============================================================
\chapter[Contexte du projet, étude de l'existant et problématique]%
{Contexte du projet, étude de l'existant\\et problématique}
\label{ch:cadre-existant}
\soustitrechapitre{De l'audit d'une plateforme en exploitation à la formulation d'un besoin de socle}

% --- Blancs autour des flottants, resserrés pour ce chapitre et le suivant
%     (valeurs LaTeX par défaut : 20 pt, 12 pt, 12 pt). Rétablis en fin de
%     chapitre~2 : le réglage ne sort pas du périmètre de cette équipe. ---
\setlength{\textfloatsep}{10pt plus 2pt minus 2pt}
\setlength{\intextsep}{8pt plus 2pt minus 2pt}
\setlength{\floatsep}{8pt plus 2pt minus 2pt}


% --- Pictogrammes de couverture (locaux à ce chapitre, réutilisés
%     à l'identique dans l'annexe A, section 1, alimentée par ce
%     chapitre) : cercle plein / mi-plein / vide, dessinés en TikZ
%     plutôt qu'en glyphes Unicode (couverture incertaine dans Times
%     New Roman) ---
% Les trois pictogrammes de couverture sont definis dans config/commands.tex :
% ils sont employes ici et dans l'annexe A, donc hors d'un fichier de chapitre.

Un projet de sécurité ne se justifie pas par les technologies qu'il emploie, mais par le
risque qu'il fait disparaître. Ce chapitre établit sur des faits vérifiés pourquoi celui-ci
devait être mené : l'entreprise d'accueil (\cref{sec:cadre-projet}), l'application qui sert de
cas réel (\cref{sec:etude-existant}), l'audit de son hébergement
(\cref{sec:critique-existant}), puis la problématique, les objectifs et la démarche
(\cref{sec:problematique-objectifs,sec:perimetre-contraintes-demarche}). L'application
fonctionne, mais repose sur un modèle périmétrique et un environnement non reproductible ; le
socle corrige ces deux défauts, et il est réutilisable parce que c'est l'intérêt de
l'entreprise.

\section{Cadre du projet et entreprise d'accueil}
\label{sec:cadre-projet}

\subsection{Fiche signalétique, activités et positionnement}
\label{subsec:entreprise-accueil}
\label{subsec:activites-positionnement}

MENAL-SARL est une société de services informatiques à deux axes --- services et fourniture de
solutions --- qui accompagne un client de la conception à l'exploitation. Le
\cref{tab:fiche-menal} en donne l'identité, l'organisation et le périmètre d'accueil du projet.

\begin{table}[!htbp]
\centering
\scriptsize
\caption{Fiche signalétique de l'entreprise d'accueil.}
\label{tab:fiche-menal}
\begin{tabular}{|P{4.7cm}|P{10.7cm}|}
\hline
\textbf{Élément} & \textbf{Information} \\
\hline
Raison sociale & MENAL-SARL \\ \hline
Forme juridique et création & Société à responsabilité limitée, créée en 2025 \\ \hline
Implantations & Siège social à Nouakchott (Mauritanie) ; implantation secondaire en France \\ \hline
Secteur et marchés & Services et technologies de l'information, pour des organisations
mauritaniennes engagées dans la digitalisation et la sécurisation de leur système
d'information \\ \hline
Axes d'activité & \textbf{Services informatiques} : architectures cloud, DevSecOps et Zero Trust, applications sécurisées, chaînes de données, audit de systèmes existants. \textbf{Fourniture de solutions IT} : matériels, installation, déploiement, maintenance \\ \hline
Organisation & Trois niveaux : direction générale, trois directions opérationnelles, neuf unités (\cref{fig:organisation-menal}) \\ \hline
Domaine institutionnel & \texttt{menal-sarl.com} \\ \hline
Périmètre d'accueil du projet & Direction technique --- unités \emph{Cloud et DevOps} et \emph{Cyber et SOC} \\ \hline
Encadrement & M.~Houssein Ezzedine (entreprise), Mme~Rihem Matoussi (académique) \\ \hline
Période du projet & Du \datedebutstage{} au \datefinstage{} \\
\hline
\end{tabular}
\end{table}

Les organisations que l'entreprise adresse sont des structures émergentes, et l'entreprise
elle-même en est une : à cette taille, ni une supervision commerciale, dont la licence est
indexée sur le volume de journaux ingérés, ni une équipe de sécurité à demeure ne s'amortissent
sur un client unique. Cette contrainte de marché est reprise telle quelle comme
\textbf{contrainte humaine} et \textbf{contrainte économique} du projet
(\cref{subsec:contraintes}), puis comme \textbf{verrou V1} de l'état de l'art
(\cref{ch:etat-art}) : le même fait à trois échelles. Le socle n'est pas la transposition d'un
modèle existant, c'est la réponse à un modèle qui ne s'applique pas.

\subsection{Organisation, gouvernance et périmètre d'accueil}
\label{subsec:service-accueil}

La \cref{fig:organisation-menal} donne l'organigramme à la date de \dateref : une direction
générale, trois directions opérationnelles et neuf unités, dont quatre à la direction technique
--- hébergement et déploiement, sécurité et supervision, développement applicatif, chaînes de
données et modèles.

\begin{figure}[!htbp]
\centering
\begin{tikzpicture}[
  x=1cm, y=0.86cm, font=\scriptsize,
  dg/.style={draw=black!80, line width=0.9pt, rounded corners=2pt, fill=black!8,
             align=center, inner sep=3pt, minimum height=0.72cm, text width=3.3cm},
  dir/.style={draw=black!75, rounded corners=2pt, fill=white,
              align=center, inner sep=3pt, minimum height=0.86cm, text width=3.4cm},
  un/.style={draw=black!60, rounded corners=2pt, fill=white,
             align=center, inner sep=2pt, minimum height=0.70cm, text width=1.44cm},
  unacc/.style={draw=black!85, line width=1.2pt, rounded corners=2pt, fill=black!10,
                align=center, inner sep=2pt, minimum height=0.70cm, text width=1.44cm},
  proj/.style={draw=black!80, rounded corners=2pt, fill=black!4,
               align=center, inner sep=3pt, minimum height=0.78cm, text width=4.6cm},
  l/.style={black!70, line width=0.6pt},
  la/.style={black!85, line width=1.0pt},
  ai/.style={-{Latex[length=1.4mm]}, black!55, densely dotted}]

% ============ niveau 1 =========================================================
\node[dg] (dg) at (8.56,7.05) {\textbf{DIRECTION GÉNÉRALE}};

% ============ niveau 2 : les trois directions ==================================
\node[dir] (dt) at (3.49,5.30) {\textbf{Direction technique}\\[1pt] \emph{CTO / Tech Lead}};
\node[dir] (dc) at (8.83,5.30) {\textbf{Direction commerciale\\et relation client}};
\node[dir] (af) at (13.36,5.30) {\textbf{Administration\\et finance} --- \emph{RAF}};

\draw[l] (dg.south) -- (8.56,6.18);
\draw[l] (3.49,6.18) -- (13.36,6.18);
\draw[l] (dt.north) -- (3.49,6.18);
\draw[l] (dc.north) -- (8.83,6.18);
\draw[l] (af.north) -- (13.36,6.18);

% ============ niveau 3 : les neuf unités =======================================
% Direction technique : quatre unités, dont deux forment le périmètre d'accueil.
\node[unacc] (u1) at (0.85,3.05) {Cloud\\et DevOps};
\node[unacc] (u2) at (2.61,3.05) {Cyber\\et SOC};
\node[un]    (u3) at (4.37,3.05) {Software\\Engi\-neering};
\node[un]    (u4) at (6.13,3.05) {Data\\et IA};
% Direction commerciale et relation client : deux unités.
\node[un]    (u5) at (7.95,3.05) {Business\\Develop\-ment};
\node[un]    (u6) at (9.71,3.05) {Account\\Manage\-ment};
% Administration et finance : trois unités.
\node[un]    (u7) at (11.60,3.05) {Finance};
\node[un]    (u8) at (13.36,3.05) {Ressources\\humaines};
\node[un]    (u9) at (15.12,3.05) {Achats};

\draw[l] (dt.south) -- (3.49,4.12);   \draw[l] (0.85,4.12) -- (6.13,4.12);
\foreach \n in {u1,u2,u3,u4}{\draw[l] (\n.north) -- (\n|-3.49,4.12);}
\draw[l] (dc.south) -- (8.83,4.12);   \draw[l] (7.95,4.12) -- (9.71,4.12);
\foreach \n in {u5,u6}{\draw[l] (\n.north) -- (\n|-8.83,4.12);}
\draw[l] (af.south) -- (13.36,4.12);  \draw[l] (11.60,4.12) -- (15.12,4.12);
\foreach \n in {u7,u8,u9}{\draw[l] (\n.north) -- (\n|-13.36,4.12);}

% ============ périmètre d'accueil : trait épais doublé par une mention =========
\draw[la] (0.09,2.30) -- (0.09,2.10) -- (3.37,2.10) -- (3.37,2.30);
\draw[la] (1.73,2.10) -- (1.73,1.90);
\node[anchor=north, font=\scriptsize\bfseries, black!85] at (1.73,1.90)
     {périmètre d'accueil du projet};

% ============ le projet ========================================================
\node[proj] (pj) at (2.40,0.62)
     {\textbf{Projet} --- socle d'hébergement sécurisé\\[1pt]
      \emph{intervenant technique unique}};
\draw[la] (1.73,1.36) -- (1.73,1.01);

% ============ chaîne de revue des livrables ====================================
\node[anchor=west, align=left, font=\scriptsize, draw=black!45, rounded corners=2pt,
      fill=black!2, inner sep=4pt, text width=8.0cm] (rev) at (7.30,0.85)
     {\textbf{Chaîne de revue des livrables} \\[1pt]
      Besoins et effets métier $\rightarrow$ responsable de l'application pilote \;·\;
      arbitrages d'architecture $\rightarrow$ encadrant entreprise \;·\;
      méthode et démonstration $\rightarrow$ encadrante académique
      \emph{(hors entreprise)}};
\draw[ai] (pj.east) -- (rev.west);

\end{tikzpicture}
\caption[Organigramme de MENAL-SARL et rattachement du projet]{Organigramme de MENAL-SARL et
rattachement du projet, à la date de \dateref. Ce qu'il faut y voir : une réalisation
mono-opérateur, portée par deux des neuf unités, validée par trois regards dont un hors de
l'entreprise.}
\label{fig:organisation-menal}
\end{figure}

Trois faits caractérisent le cadre d'exécution, et chacun est repris tel quel comme donnée
d'entrée dans la suite du mémoire. Le \textbf{périmètre d'accueil} joint deux unités,
\emph{Cloud et DevOps} et \emph{Cyber et SOC}, et ce rattachement double n'est pas une commodité
d'organigramme : le socle est précisément l'objet qui les joint --- de la première seule
relèverait une plateforme d'hébergement, de la seconde seule un outillage de supervision posé sur
une infrastructure qu'il ne maîtrise pas. Le projet est ensuite \textbf{mono-porteur} ---
l'élève-ingénieur en assure seul la conception, la réalisation et l'exploitation ---, donnée
reprise comme contrainte humaine en \cref{subsec:contraintes}, qui écarte par construction tout
composant à exploitation permanente. Les livrables suivent enfin une \textbf{chaîne de revue à
trois niveaux} dont le troisième regard est extérieur à l'entreprise : c'est ce qui compense la
faiblesse structurelle d'un contrôle mono-opérateur.

\subsection{Mission confiée et positionnement stratégique}
\label{subsec:mission-confiee}

La mission comporte cinq volets : auditer l'application pilote, formaliser les risques qu'elle
porte, concevoir un socle Zero Trust et DevSecOps, l'implémenter en infrastructure décrite en code
et confronter ses contrôles à des tests et à des mesures. J'en ai assuré la continuité de bout en
bout ; le planning figure en \cref{sec:planning-travail}.

Le point décisif est ailleurs, et il commande l'ampleur du travail. Un prestataire qui héberge
les applications de ses clients dispose de deux voies. La première construit un environnement sur
mesure par client : rapide au premier contrat, elle fait recommencer à chaque suivant le réseau,
les identités, les secrets, la journalisation et la chaîne de livraison, et le niveau de sécurité
atteint est celui qu'autorise le budget du client le plus modeste. La seconde construit une fois
un socle standardisé dont chaque client devient un locataire : le coût de conception est payé une
fois, le niveau de sécurité est celui du socle et non du contrat, et l'accueil d'une application
se réduit à l'instanciation de paramètres. \textbf{Le coût de sécurité cesse alors d'être une
dépense récurrente pour devenir un actif de l'entreprise} : c'est la voie retenue, et elle
justifie qu'un projet de fin d'études porte sur un socle plutôt que sur la sécurisation d'une
application. Elle devient une exigence vérifiable au \cref{ch:besoins-menaces} --- le besoin
\textbf{BF13} interdit à tout composant du socle une logique propre à une application ---,
confrontée à sa preuve au \cref{ch:realisation} par l'accueil d'un second locataire puis
discutée au \cref{ch:validation}.

\section{L'application pilote}
\label{sec:etude-existant}

\subsection{Nature, publics et modèle fonctionnel}
\label{subsec:appli-pilote}

L'application retenue comme premier cas d'usage du socle est ELSON, une plateforme de production
participative gamifiée qui construit des jeux de données parallèles, texte et audio, pour les
langues nationales de Mauritanie, en commençant par l'arabe hassaniya. Son public est constitué
de locuteurs hassanophones de Mauritanie et de la diaspora, mobilisés par une compétition dotée
de prix, et ses données de meilleure qualité alimentent des modèles de reconnaissance vocale, de
traduction et de synthèse. Le cycle fonctionnel est une boucle à trois temps
(\cref{fig:elson-fonctionnel}) : un contributeur traduit des phrases sources, enregistre sa voix,
puis évalue les contributions d'autrui ; les évaluations font le classement, le classement
commande la dotation, et les contributions validées alimentent le corpus. Chacun de ces trois
temps produit une donnée d'un genre différent et offre une prise différente à un adversaire.

\begin{figure}[!htbp]
\centering
\begin{tikzpicture}[
  x=1cm, y=0.88cm, font=\scriptsize,
  temps/.style={draw=black!70, fill=black!4, align=center,
                inner sep=3pt, minimum height=0.72cm},
  actif/.style={draw=black!70, double, double distance=0.7pt, fill=white,
                align=center, inner sep=3pt, minimum height=0.72cm},
  ext/.style={draw=black!70, fill=black!10, align=center,
              inner sep=3pt, minimum height=0.55cm},
  cadre/.style={draw=black!55, densely dotted, rounded corners=4pt, line width=0.7pt},
  a/.style={-{Latex[length=1.5mm]}, black!65},
  ai/.style={-{Latex[length=1.5mm]}, black!55, densely dotted},
  et/.style={font=\scriptsize, fill=white, inner sep=1.5pt}]

% ============ cadre englobant les trois temps de la boucle ====================
\draw[cadre] (4.10,1.25) rectangle (9.35,7.10);
\node[anchor=south, align=center, font=\scriptsize\itshape, text width=5.4cm, black!60]
      at (6.72,7.18) {application web progressive installable\\--- aucun canal natif};

% ============ point d'entree unique ===========================================
\node[ext, text width=1.1cm] (net) at (1.00,4.20) {Internet};
\draw[a, line width=0.9pt] (net.east) -- (4.10,4.20);
\node[anchor=south, align=center, text width=2.1cm, font=\scriptsize, inner sep=1pt]
      at (3.00,4.32) {HTTPS --- point d'entrée unique};

% ============ les trois temps de la boucle, sens horaire ======================
\node[temps, text width=2.2cm] (tr) at (6.05,6.30) {\textbf{Traduire}};
\node[temps, text width=1.8cm] (en) at (8.15,4.20) {\textbf{Enregistrer}\\\textbf{sa voix}};
\node[temps, text width=2.6cm] (ev) at (6.15,2.20)
      {\textbf{Évaluer} les\\contributions d'autrui};

\draw[a] (tr.south east) -- node[et, pos=0.50] {phrase source}      (en.north);
\draw[a] (en.south west) -- node[et, pos=0.50] {contribution audio} (ev.north);
\draw[a] (ev.west) -- (4.55,2.20) -- (4.55,6.30) -- (tr.west);
\node[rotate=90, et, font=\scriptsize] at (4.55,5.35) {note de validation};

% ============ les trois actifs produits =======================================
\node[actif, text width=4.2cm] (cor) at (12.85,6.50) {Corpus texte et voix};
\node[actif, text width=4.2cm] (pro) at (12.85,4.20)
      {Profil du contributeur :\\identité, voix, empreinte};
\node[actif, text width=4.2cm] (cla) at (12.85,1.90) {Classement et dotation};

% -- alimentation du corpus : le texte puis la voix
\draw[a] (tr.north east) -- (cor.west);
\draw[a] (en.north east) -- (cor.south west);
% -- alimentation du classement
\draw[a] (ev.east) -- (cla.west);
% -- alimentation du profil, par les trois temps
\draw[ai] (tr.east)       -- (pro.north west);
\draw[ai] (en.east)       -- (pro.west);
\draw[ai] (ev.north east) -- (pro.south west);

\end{tikzpicture}
\caption[Boucle fonctionnelle de l'application pilote et actifs produits]{Boucle fonctionnelle
de l'application pilote et actifs produits. Chacun des trois temps produit un actif d'une nature
différente --- corpus, classement doté, profil de contributeur --- alors que tous trois
transitent par un seul point d'entrée HTTP.}
\label{fig:elson-fonctionnel}
\end{figure}

La \cref{fig:capture-K35} restitue le premier de ces trois temps tel qu'il se présente au
contributeur.

\begin{figure}[!htbp]
\centering
\IfFileExists{figures/K35.png}{%
  \setlength{\fboxrule}{0.4pt}\setlength{\fboxsep}{0pt}%
  \fcolorbox{black!30}{white}{\includegraphics[width=0.62\textwidth]{figures/K35.png}}%
}{%
  \begin{minipage}{0.62\textwidth}
  \begin{extraitconf}[breakable=false, fontupper=\footnotesize,
      title={Cadre de capture K35 : spécification de prise}]
  \textbf{Vue :} écran de contribution de l'application pilote, environnement de recette --- la vue qu'obtient un contributeur au premier temps de la boucle : phrase source affichée, langue sélectionnée, commande d'enregistrement de la voix.\par\smallskip
  \textbf{Date de prise :} sur la révision de recette en service. \emph{La date portée par la capture fait foi.}\par\smallskip
  \textbf{Contenu probant :} la boucle décrite ci-dessus servie par le système réel, et non par une maquette : une phrase source du corpus, la langue de contribution et la commande d'enregistrement visibles dans le même cadre. \textbf{Aucune contribution réelle d'un tiers ne doit apparaître} --- ni voix, ni texte soumis, ni identité de contributeur : la vue est prise sur un compte de recette.\par\smallskip
  \textbf{Masquage} (rectangle opaque, jamais de flou) \textbf{:} nom et identifiant du compte connecté, adresse du service.
  \end{extraitconf}
  \end{minipage}%
}
\caption[Écran de contribution de l'application pilote (environnement de recette)]{Écran de contribution de l'application pilote, environnement de recette : le premier temps de la boucle de la \cref{fig:elson-fonctionnel}. Restitution d'état ; la date de relevé est portée par la capture.}
\label{fig:capture-K35}
\end{figure}

Une précision d'architecture porte ici un argument et non un détail. ELSON est une
\textbf{application web progressive installable} servie par un frontend à rendu serveur --- ni
application native, ni application hybride, le caractère installable tenant à un manifeste et à
un agent de service écrits à la main. Toute la surface d'attaque transite donc par HTTP : un
\textbf{filtrage applicatif} placé au bord la couvre entièrement, et \emph{aucun canal mobile
n'échappe au périmètre du socle}.

Il ne s'agit pas d'un prototype mais d'une application en exploitation réelle. À la date de
l'audit du 29/07/2026, le dépôt comptait environ 34\,600~lignes de code TypeScript et
225~points d'entrée HTTP répartis en onze routeurs, dont environ 165~derrière un contrôle
d'administration ; la base réunissait 54~tables PostgreSQL, 17~vues, 15~fonctions et
déclencheurs, définies dans 74~fichiers SQL dont 72~migrations. Ces valeurs sont datées parce
qu'elles décrivent un logiciel vivant --- au 28/08/2026, environ 245~points d'entrée et
85~migrations ---, et l'analyse repose sur l'état du 29/07/2026, celui sur lequel les décisions
ont été prises.

\subsection{Sensibilité des données traitées}
\label{subsec:donnees-sensibles}

Le niveau de sécurité exigé d'un hébergement ne se déduit pas du nombre de ses utilisateurs mais
de la nature de ce qu'il conserve. Le \cref{tab:sensibilite-elson} recense les catégories
manipulées, désignées par le champ qui les porte, et l'effet de leur exposition.

{\setlength{\tabcolsep}{3pt}
\setlength{\LTpre}{4pt}\setlength{\LTpost}{4pt}
\begin{longtable}{|P{3.3cm}|P{4.3cm}|P{7.7cm}|}
\caption{Catégories de données traitées par l'application pilote et effet d'une exposition.}
\label{tab:sensibilite-elson}\\
\hline
\textbf{Champ ou groupe de champs} & \textbf{Nature au sens de la protection des données} &
  \textbf{Ce que produit son exposition} \\
\hline
\endfirsthead
\hline
\textbf{Champ ou groupe de champs} & \textbf{Nature au sens de la protection des données} &
  \textbf{Ce que produit son exposition} \\
\hline
\endhead
\hline
\endfoot
\texttt{nni} & Identifiant national officiel, unique par compte et rendu immuable par
déclencheur & Usurpation d'identité administrative hors de la plateforme ; l'identifiant ne
peut pas être réémis \\ \hline
\texttt{audio\_url} et fichiers associés & Enregistrement vocal du contributeur, donc donnée
biométrique au sens de l'article~9 du RGPD & Identification et imitation vocale ; la voix ne
peut pas être changée \\ \hline
\texttt{whatsapp}, \texttt{first\_name}, \texttt{last\_name}, \texttt{birthdate},
\texttt{email} & Contact direct et état civil & Joindre la personne hors de la plateforme et
rapprocher son compte de son identifiant national \\ \hline
\texttt{signup\_ip}, \texttt{signup\_geo}, \texttt{signup\_ua},
\texttt{device\_fingerprint} & Adresse de connexion, géolocalisation et empreinte d'appareil,
indexées pour le regroupement anti-fraude & Relie des comptes distincts et localise une
personne : l'outil anti-fraude devient un outil de surveillance s'il fuit \\ \hline
\texttt{fraud\_score} & Décision individuelle automatisée : au-delà d'un seuil, l'inscription
est rejetée sans intervention humaine & Relève de l'article~22 du RGPD : la décision doit
pouvoir être expliquée et contestée, donc son historique protégé \\ \hline
\texttt{contributor\_credits} (nom et numéro d'un tiers) & Données personnelles de personnes
\emph{non inscrites} sur la plateforme & Le consentement est saisi par le contributeur, pas
par la personne concernée : l'exposition atteint des tiers sans relation avec le service
\\ \hline
\texttt{audit\_log.details} & Traces d'administration portant des données personnelles en
clair & Étend la portée de toutes les lignes précédentes à la journalisation elle-même \\
\end{longtable}
}

\textbf{L'irréversibilité d'abord.} L'identifiant national et l'empreinte vocale sont les deux
seules catégories qui ne se révoquent pas : leur exposition est définitive et le reste après la
fin de vie du service. Au 28/08/2026, elles sont protégées par le chiffrement d'infrastructure du
socle et, pour l'audio, par des adresses signées à durée courte, sans chiffrement applicatif
propre : c'est ce constat qui fixe le niveau d'exigence attendu de l'hébergement.

\textbf{L'incitation ensuite.} Aucun flux de paiement ne transite par l'application, ce qui
restreint utilement le périmètre de conformité ; mais la compétition est dotée de prix, et cette
incitation crée un adversaire que la plupart des modèles de menace oublient : celui qui
\emph{dispose déjà d'un accès légitime}, aux requêtes bien formées et aux quotas respectés ---
le \textbf{participant fraudeur} (acteur A9 du \cref{ch:besoins-menaces}). C'est la
justification la plus directe du modèle Zero Trust : un périmètre ne dit rien d'un adversaire
déjà admis.

\textbf{Le cadre juridique enfin.} Le traitement s'exécute dans une région européenne et
concerne des personnes situées en Mauritanie et dans la diaspora : le RGPD s'applique par le lieu
du traitement, et pour la diaspora au titre de son article~3. Trois dispositions se cumulent ---
articles~9 (biométrie), 22 (décision automatisée) et 35 ---, et ce cumul déclenche l'obligation
d'analyse d'impact. Le rattachement institutionnel du projet appelle en outre la \textbf{loi
tunisienne n\textsuperscript{o}~2004-63 du 27~juillet 2004} : formalité préalable, consentement
écrit, droits d'accès et de rectification, encadrement des transferts hors territoire. Reste le
droit mauritanien, alors que les personnes concernées sont mauritaniennes et que l'identifiant
collecté est un identifiant d'État : aucune norme mauritanienne n'ayant été vérifiée dans ce
travail, le mémoire pose la question comme \textbf{ouverte et assumée} plutôt que de la refermer
par une référence approximative. Son traitement relève du responsable de traitement
(\cref{subsec:perimetre}).

\subsection{Trois actifs, trois adversaires}
\label{subsec:actifs-adversaires}

L'analyse de l'application ne fait pas apparaître un actif à protéger, mais trois, de natures
différentes et exposés à des adversaires différents. \textbf{Le corpus} est le résultat
économique du projet --- données texte et voix validées, constituées au prix d'un effort humain
long, exportables par un point d'entrée dédié --- et il attire l'\textbf{acteur intéressé par le
corpus}, patient et ciblant délibérément (A8). \textbf{L'équité de la compétition} tient au
classement, qui commande la dotation : faussé, il ne détourne pas seulement un prix, il pollue le
corpus qu'il devait sélectionner ; l'adversaire y est le \textbf{participant fraudeur},
disposant d'un accès légitime et d'une incitation financière (A9). \textbf{L'identité des
contributeurs} --- voix et identifiant national --- est la seule catégorie dont l'exposition ne
se répare pas, et elle expose à l'\textbf{attaquant opportuniste} (A7) comme à
l'\textbf{interne ou prestataire} disposant d'un accès légitime initial (A11).

C'est cette mise en regard, et non une liste de bonnes pratiques, qui justifie le niveau de
défense construit ensuite : contrôler ce qui entre (l'opportuniste), contraindre ce qui est déjà
admis (le fraudeur), prouver ce qui s'est passé (l'interne). Le \cref{ch:besoins-menaces} reprend
ces trois actifs comme valeurs métier de l'analyse de risque.

\subsection{L'hébergement initial et son modèle de sécurité}
\label{subsec:hebergement-initial}

Avant le projet, l'application est déployée sur un serveur unique loué chez un hébergeur
européen : des conteneurs orchestrés par un seul fichier de composition --- frontend Next.js,
API Express en cluster, base PostgreSQL~17, pooler PgBouncer, proxy Caddy (TLS), passerelle
WhatsApp tierce et sauvegarde locale. Ce modèle (\cref{fig:archi-avant}) repose sur un pare-feu
unique en entrée et une confiance implicite à l'intérieur.

\begin{figure}[!htbp]
\centering
\begin{tikzpicture}[
  y=0.9cm,
  composant/.style={draw=black!70, thick, rounded corners=2pt, fill=blue!3,
    minimum width=1.85cm, minimum height=0.9cm, align=center, font=\scriptsize},
  externe/.style={composant, fill=white, minimum width=1.3cm},
  critique/.style={composant, draw=red!70!black, very thick, fill=red!5},
  flux/.style={-{Latex[length=1.8mm]}, thick, draw=blue!55!black},
  retour/.style={{Latex[length=1.6mm]}-{Latex[length=1.6mm]}, thick, draw=blue!55!black},
  etiquette/.style={font=\scriptsize, inner sep=1pt}
]
\node[externe] (net) at (-1.00,0) {Internet};
\node[composant] (proxy) at (2.25,0) {Caddy\\terminaison TLS};
\node[composant] (front) at (5.00,0) {Frontend\\Next.js};
\node[composant] (api) at (7.60,0) {API Express\\4 processus};
\node[critique] (pgb) at (10.25,0) {PgBouncer\\sans\\authentification};
\node[composant] (pg) at (12.85,0) {PostgreSQL 17};
\node[composant] (waha) at (7.60,-1.75) {Passerelle WAHA\\service tiers};
\node[composant] (backup) at (12.85,1.75) {Sauvegarde\\locale};

\begin{scope}[on background layer]
  \node[draw=black!60, dashed, thick, rounded corners=4pt, fill=black!4,
    fit={(proxy)(front)(api)(pgb)(pg)(waha)(backup)}, inner sep=8pt,
    label={[font=\footnotesize\bfseries]above:Serveur unique --- réseau interne partagé}] {};
\end{scope}

\draw[flux] (net) -- (proxy) node[etiquette, pos=0.35, above] {HTTPS};
\draw[flux] (proxy) -- (front) node[etiquette, midway, above] {HTTP};
\draw[flux] (front) -- (api) node[etiquette, midway, above] {API};
\draw[flux] (api) -- (pgb) node[etiquette, midway, above] {SQL};
\draw[flux] (pgb) -- (pg) node[etiquette, midway, above] {SQL};
\draw[retour] (api.south) -- (waha.north) node[etiquette, midway, right] {messages};
\draw[flux] (pg.north) -- (backup.south) node[etiquette, midway, right] {copie};

\node[draw=red!70!black, rounded corners=2pt, fill=red!5, font=\scriptsize,
  align=center, text width=3.0cm] (b3) at (11.05,-1.75)
  {\textbf{B3} --- PgBouncer accepte tout conteneur interne};
\draw[-{Latex[length=1.8mm]}, thick, draw=red!70!black]
  (b3.north) -- (pgb.south);
\end{tikzpicture}
\caption[Modèle de sécurité de l'hébergement initial]{Modèle de sécurité avant le projet :
périmètre unique en entrée, confiance implicite à l'intérieur. Tout composant présent sur le
réseau interne est traité comme légitime --- c'est la carence C1.}
\label{fig:archi-avant}
\end{figure}

\section{Audit de l'existant et critique}
\label{sec:critique-existant}

\subsection{Méthode et résultats de l'audit}
\label{subsec:methode-audit}
\label{subsec:resultats-audit}

Un audit technique et sécurité a été réalisé le 29/07/2026, en amont de toute conception,
par analyse statique exhaustive du dépôt --- code, configurations, scripts de déploiement et
74~fichiers SQL ---, sans exécution ni déploiement. Deux règles l'ont encadré, qui gouvernent
aussi l'écriture de ce mémoire : la \textbf{règle de preuve} (seul ce qui est dans le dépôt est
documenté, chemin et ligne cités) et la \textbf{règle de confidentialité} (aucun secret
reproduit, valeurs sensibles désignées par leur nom de variable). Chaque constat
porte l'un de quatre statuts --- \emph{vérifié}, \emph{dépendant de l'exécution} donc non
observable dans le dépôt seul, \emph{absent}, \emph{recommandé} ; la répartition des
41~critères par catégorie et par statut est donnée par la \cref{graph:audit-41}.

L'audit ne conclut pas à un travail de mauvaise qualité : plusieurs mécanismes dépassent ce que
l'on observe à cette taille --- aucune injection SQL ou système, sessions avec rotation et
détection de réutilisation, intégrité de la compétition assurée par quatorze signaux de fraude.
La logique applicative est mûre ; le problème est dans la configuration et l'industrialisation.
Au 29/07/2026, l'audit retient six constats bloquants.

\begin{itemize}\setlength{\itemsep}{1pt}
  \item \textbf{B1} --- compte administrateur par défaut, au mot de passe en clair dans le schéma versionné ;
  \item \textbf{B2} --- secret unique partagé par sept mécanismes cryptographiques ;
  \item \textbf{B3} --- pooler de base de données accessible sans authentification ;
  \item \textbf{B4} --- jeton de téléchargement du corpus non lié à un utilisateur, à une adresse IP ou à un usage unique ;
  \item \textbf{B5} --- schéma de base non reconstructible depuis le dépôt, avec trois objets sans définition versionnée ;
  \item \textbf{B6} --- état applicatif local au processus, incompatible avec la mise à l'échelle horizontale.
\end{itemize}

Le mémoire conserve ces six constats dans l'état daté où ils ont été relevés, parce que ce sont
eux qui fondent le besoin de socle ; plusieurs ont depuis été traités par l'éditeur, en partie à
l'occasion de l'accueil de l'application sur le socle (\cref{ch:realisation}). La
\cref{graph:audit-41} donne l'écart à lire : l'application fait ce qu'elle doit faire, mais elle
n'est ni reproductible, ni observable, ni conforme, ni déployable de façon maîtrisée.

\begin{figure}[!htbp]
\centering
\begin{tikzpicture}[y=0.28cm, x=1cm, font=\footnotesize]
\def\ymax{18}
\draw[thick] (-0.4,0) -- (11.6,0);
\draw[thick] (-0.4,0) -- (-0.4,\ymax);
\node[rotate=90, anchor=south, font=\scriptsize] at (-1.35,9) {Nombre de critères};
\foreach \y in {0,5,10,15}{
  \draw (-0.55,\y) -- (-0.4,\y);
  \node[anchor=east, font=\scriptsize] at (-0.6,\y) {\y};
  \draw[black!20] (-0.4,\y) -- (11.6,\y);
}
\foreach \cat/\x/\a/\b/\c/\t in {%
  Fonctionnel/0.4/6/1/0/7, Technique/2.9/0/1/9/10, Sécurité/5.4/6/4/7/17, Conformité/7.9/0/1/6/7}{
  \fill[black!8]  (\x,0) rectangle ++(1.6,\a);
  \draw[thick]    (\x,0) rectangle ++(1.6,\a);
  \fill[black!30] (\x,\a) rectangle ++(1.6,\b);
  \draw[thick]    (\x,\a) rectangle ++(1.6,\b);
  \fill[black!65] (\x,{\a+\b}) rectangle ++(1.6,\c);
  \draw[thick]    (\x,{\a+\b}) rectangle ++(1.6,\c);
  \node[anchor=north, font=\scriptsize] at ({\x+0.8},-0.3) {\cat};
  \node[anchor=south, font=\scriptsize\bfseries] at ({\x+0.8},\t) {\t};
}
\fill[black!8]  (10.0,15.4) rectangle ++(0.45,1.1); \draw[thick] (10.0,15.4) rectangle ++(0.45,1.1);
\node[anchor=west, font=\scriptsize] at (10.55,15.95) {Prêt (12)};
\fill[black!30] (10.0,13.6) rectangle ++(0.45,1.1); \draw[thick] (10.0,13.6) rectangle ++(0.45,1.1);
\node[anchor=west, font=\scriptsize] at (10.55,14.15) {Partiel (7)};
\fill[black!65] (10.0,11.8) rectangle ++(0.45,1.1); \draw[thick] (10.0,11.8) rectangle ++(0.45,1.1);
\node[anchor=west, font=\scriptsize] at (10.55,12.35) {Non prêt (22)};
\node[anchor=north west, font=\scriptsize, align=left] at (-0.4,-1.6)
  {Chaque barre est empilée dans l'ordre prêt / partiel / non prêt ; le total figure au sommet.};
\end{tikzpicture}
\caption[Répartition des 41~critères de préparation au déploiement]{Répartition des
41~critères de préparation au déploiement, par catégorie et par statut (audit du 29/07/2026).}
\label{graph:audit-41}
\end{figure}

\subsection{Cinq carences structurantes et positionnement du socle}
\label{subsec:carences-structurantes}
\label{subsec:distinction-appli-socle}
\label{subsec:positionnement-modeles}

Reformulés à un niveau d'abstraction plus élevé, ces constats font apparaître cinq carences
structurantes qui ne sont pas propres à cette application : elles caractérisent un système
construit sans socle.

\begin{itemize}\setlength{\itemsep}{1pt}
  \item \textbf{C1 --- Modèle de confiance périmétrique} : barrière d'entrée unique et confiance implicite à l'intérieur ;
  \item \textbf{C2 --- Environnement non reproductible} : infrastructure et schéma de données non décrits en code ;
  \item \textbf{C3 --- Chaîne de livraison non contrôlée} : aucun contrôle automatique entre l'écriture du code et sa mise en production ;
  \item \textbf{C4 --- Absence d'observabilité et de détection} : événements non centralisés, non corrélés, intrusion non détectée ;
  \item \textbf{C5 --- Absence de cadre de conformité} : données d'identité et biométriques traitées sans base légale, sans durée de conservation ni droit d'effacement.
\end{itemize}

Une conclusion importante en découle : corriger l'application ne suffit pas. Les carences C1 à C5
ne sont pas des défauts de code mais des propriétés absentes de l'environnement d'exécution ;
aucun correctif applicatif ne créera d'identité par charge de travail, de journalisation
centralisée, de chaîne de livraison contrôlée ni de reproductibilité. D'où le périmètre : le code
applicatif en reste dehors, le socle qui héberge, isole, contrôle, journalise et supervise
l'application en est l'objet --- et un socle conçu indépendamment d'elle peut en accueillir
d'autres (\cref{subsec:mission-confiee}).

Trois modèles situent la solution visée : le modèle~A, l'architecture périmétrique de l'existant
(\cref{subsec:hebergement-initial}) ; le modèle~B, un hébergement cloud managé sans démarche de
sécurité explicite ; le modèle~C, le socle Zero Trust visé. La comparaison à quinze critères,
issus pour la plupart des principes Zero Trust~\cite{nist800207}, figure au
\cref{tab:positionnement-15-complet}. Le modèle~C domine sur la sécurité et la reproductibilité
et perd sur trois critères, repris au \cref{ch:validation} : la simplicité initiale, la mise à
l'échelle horizontale, écartée des objectifs (\cref{subsec:perimetre}), et la maturité
opérationnelle, une supervision sur mesure ne bénéficiant d'aucun support éditeur. Ce sont ces
trois pertes qui rendent crédible le reste de la comparaison.

\section{Problématique et objectifs}
\label{sec:problematique-objectifs}

\subsection{Problématique et questions dérivées}
\label{subsec:problematique}

Les constats précédents permettent de formuler la question centrale du mémoire.
\begin{keybox}
Comment concevoir, déployer et surtout \emph{prouver} un socle d'hébergement cloud aligné sur
les principes Zero Trust et sur une chaîne de livraison DevSecOps ? La question se pose dans
les conditions réelles d'une entreprise émergente : une seule personne pour l'exploitation,
un budget de quelques dizaines d'euros par mois, aucune solution de supervision commerciale.
Et quel prix faut-il payer, en faux positifs, en latence de détection et en complexité, pour
atteindre ce résultat ?
\end{keybox}

Le verbe \emph{prouver} est central : chaque contrôle doit être associé à un test qui démontre
son efficacité, non à une déclaration de conformité ; et la question porte un coût explicite, une
architecture qui ne mesure pas ce qu'elle coûte n'étant pas une architecture d'ingénieur. Six
questions dérivées en découlent, chacune associée à un objectif et à un chapitre. \textbf{Q1} (O2 · ch.~5-6) un environnement complet peut-il être
détruit puis reconstruit à l'identique depuis le seul dépôt, en combien de temps ? \textbf{Q2}
(O1 · ch.~4-6) que reste-t-il atteignable depuis Internet quand le point d'entrée unique est
contourné ? \textbf{Q3} (O1 · ch.~4-6) que peut faire un compte ou un service compromis,
jusqu'où peut-il aller ? \textbf{Q4} (O4 · ch.~5-6) quelles attaques sont détectées, lesquelles
ne le sont pas, en combien de temps ? \textbf{Q5} (O3 · ch.~5-6) qu'est-ce qui empêche une
régression de sécurité d'atteindre la production, à quel coût en blocages injustifiés ?
\textbf{Q6} (O5 · ch.~5-6) les vulnérabilités livrées peuvent-elles être priorisées selon les
menaces réellement observées ? Toutes portent sur des capacités, non sur des produits ; Q6 est
l'axe original du travail, reliant chaîne de développement et détection d'incidents.

\subsection{Objectifs et critères de succès}
\label{subsec:objectifs}

Un objectif qui ne peut pas être mesuré ne peut pas être défendu : chacun porte un indicateur,
une cible et un moyen de vérification (\cref{tab:objectifs}) --- contrat d'évaluation du mémoire,
auquel le \cref{ch:validation} répond point par point.

{\setlength{\tabcolsep}{3pt}
\setlength{\LTpre}{4pt}\setlength{\LTpost}{4pt}
\begin{longtable}{|P{0.7cm}|P{3.7cm}|P{5.2cm}|P{5.4cm}|}
\caption{Objectifs et critères de succès du projet.}
\label{tab:objectifs}\\
\hline
\textbf{\#} & \textbf{Objectif} & \textbf{Indicateur} & \textbf{Cible et vérification} \\
\hline
\endfirsthead
\hline
\textbf{\#} & \textbf{Objectif} & \textbf{Indicateur} & \textbf{Cible et vérification} \\
\hline
\endhead
\hline
\endfoot
O1 & Supprimer la confiance implicite du réseau & Nombre de tentatives de contournement
bloquées lors de la campagne de tests adverses & 100~\% des scénarios définis ; campagne de
tests adverses, \cref{ch:validation} \\ \hline
O2 & Rendre l'infrastructure reproductible & Part du socle décrite en code d'infrastructure ;
durée de reconstruction complète d'un environnement & 100~\% du socle ; durée mesurée et
documentée par destruction puis reconstruction chronométrée \\ \hline
O3 & Contrôler la chaîne de livraison & Nombre de portes bloquantes ; nombre de clés de
compte de service exportées & 3~portes (secrets, analyse statique, vulnérabilités) ; 0~clé ;
journaux d'exécution de la chaîne de livraison applicative \\ \hline
O4 & Détecter les événements de sécurité & Nombre de règles de détection ; couverture des
tactiques de la matrice de référence ; taux de fausses alertes sur 7~jours & Couverture
mesurée et publiée, lacunes comprises ; requêtes de détection et mesures, \cref{ch:validation} \\ \hline
O5 & Enrichir les alertes de façon déterministe & Précision de l'association alerte~$\to$~technique
d'attaque, comparée à une méthode de référence simple & Gain mesuré, positif ou négatif, sur
un jeu annoté ; évaluation comparative, \cref{ch:validation} \\ \hline
O6 & Rester exploitable et économiquement soutenable & Coût mensuel réel ; disponibilité et
latence observées ; existence de procédures d'incident exécutées au moins une fois & Quelques
dizaines d'euros par mois ; objectifs de service instrumentés ; facturation, supervision et
exécution réelle des procédures \\
\end{longtable}
}

Deux précisions : l'objectif O4 demande de publier ce que le système \textbf{ne détecte pas}, et
l'objectif O5 admet par avance un résultat négatif, condition d'une évaluation rigoureuse.

\section{Périmètre, contraintes et démarche}
\label{sec:perimetre-contraintes-demarche}

\subsection{Périmètre et hypothèses assumées}
\label{subsec:perimetre}

Sont dans le périmètre : le socle d'hébergement cloud (filtrage en entrée, identité, exécution
des charges, réseau, données, enrichissement, observabilité), la chaîne de livraison applicative,
la supervision de sécurité sur un entrepôt de données, l'enrichissement sémantique des alertes,
le tableau de bord de l'analyste et des recommandations de conformité. Restent hors périmètre :
le code interne de l'application hébergée, le déploiement multi-région et la reprise d'activité
inter-région, la mise à l'échelle massive. Le socle prend en charge trois obligations attachées
aux données de \cref{subsec:donnees-sensibles} : région unique et documentée, chiffrement par
clés gérées des secrets et de l'espace média, journalisation centralisée des accès
d'administration. Base légale, durée de conservation et droit d'effacement relèvent du
responsable de traitement : c'est la carence C5, formulée ici et non refermée par ce projet.

Trois hypothèses sont assumées. Les deux environnements réels et la production cible sont des
projets cloud distincts d'une même région, bâtis sur les mêmes modules, seules les variables de
dimensionnement changeant : \textbf{la sécurité n'est jamais un paramètre d'environnement}. La
mise à l'échelle est écartée des objectifs, ses limites étant mesurées plutôt que contournées. La
simplicité est érigée en critère d'architecture : tout composant justifie son existence, et tout
retrait se documente au même titre qu'un ajout.

\subsection{Contraintes du projet}
\label{subsec:contraintes}

Six contraintes structurent la conception ; les deux premières traduisent, à l'échelle du projet,
la contrainte de marché de \cref{subsec:activites-positionnement}.

\begin{itemize}\setlength{\itemsep}{1pt}
  \item \textbf{Humaine} : une seule personne conçoit, réalise et exploite le socle --- écarte les composants à exploitation permanente, impose des procédures écrites ;
  \item \textbf{Économique} : quelques dizaines d'euros par mois --- écarte les solutions de supervision commerciales, impose un filtrage des journaux à la source ;
  \item \textbf{Technique, exécution} : l'application hébergée écrit sur le disque local et conserve son état en mémoire, ce qui conditionne le mode d'exécution retenu ;
  \item \textbf{Technique, données} : le schéma de base n'est pas reconstructible depuis le dépôt, ce qui rend sa consolidation préalable à toute migration ;
  \item \textbf{Réglementaire} : des données d'identité officielle et biométriques imposent une région d'hébergement documentée et un chiffrement par clés gérées ;
  \item \textbf{Temporelle} : la durée du projet de fin d'études impose le phasage à critères de sortie présenté en \cref{subsec:demarche}.
\end{itemize}

\subsection{Démarche méthodologique}
\label{subsec:demarche}

La démarche suit une règle unique : rien n'existe dans l'infrastructure qui n'existe d'abord sous
forme de code versionné, les commandes interactives étant réservées aux opérations ponctuelles
--- rotation d'un secret, réponse à un incident --- et jamais au provisionnement. Elle répond à
la carence C2 : reproductibilité complète, contrôle de la sécurité avant création, détection de
toute dérive. Deux registres accompagnent le projet, tous deux tenus au \cref{ch:conception} :
celui des décisions d'architecture et celui des écarts entre conception cible et implémentation
réelle. Le projet est organisé en huit phases, chacune fermée par un critère de sortie
vérifiable, qui sont le découpage du planning (\cref{sec:planning-travail}).

\begin{itemize}\setlength{\itemsep}{1pt}
  \item \textbf{PH0 --- Immersion et cadrage} ; sortie : périmètre et interlocuteurs établis ;
  \item \textbf{PH1 --- Étude de l'existant}, par analyse statique du dépôt et de l'hébergement en place ; sortie : audit figé sur 41~critères et six constats bloquants ;
  \item \textbf{PH2 --- Risques et exigences} : valeurs métier, sources de risque, scénarios, vingt exigences ; sortie : matrice menace --- exigence --- contrôle --- test complète ;
  \item \textbf{PH3 --- Architecture Zero Trust} : vues d'architecture, sept flux, registre des décisions ; sortie : conception justifiée par écrit, chaque exigence rattachée à un contrôle ;
  \item \textbf{PH4 --- Infrastructure en code et chaînes de livraison} : réseau, identités, données, portes de contrôle, identité fédérée ; sortie : infrastructure entièrement décrite en code et planification sans différence ;
  \item \textbf{PH5 --- Détection et enrichissement} : collecte, normalisation, sept règles, couche sémantique ; sortie : chaîne de bout en bout fonctionnelle et détections visibles dans le tableau de bord ;
  \item \textbf{PH6 --- Validation et corrections} : tests adverses, mesures de détection, de performance, de résilience et de coût ; sortie : chaque objectif O1 à O6 confronté à une mesure ;
  \item \textbf{PH7 --- Rapport, transfert et clôture} ; sortie : mémoire remis et socle exploitable par un tiers à partir du seul dépôt.
\end{itemize}

\section*{Conclusion du chapitre}
\addcontentsline{toc}{section}{Conclusion du chapitre}

\looseness=-1
Ce chapitre a établi le point de départ sur des faits vérifiés. Sur un marché où ni une supervision
commerciale ni une équipe de sécurité dédiée ne s'amortissent sur un client unique, le socle
réutilisable est une décision d'entreprise et non un raffinement technique. L'application retenue
traite un identifiant national, des enregistrements vocaux et une décision automatisée --- trois
actifs, trois adversaires ; mature fonctionnellement, 6~critères sur 7, elle laisse 22 des
41~critères de préparation au déploiement non remplis au 29/07/2026, dont six points bloquants qui
se ramènent à cinq carences structurantes, toutes situées dans l'environnement d'exécution. L'objet
du mémoire n'est donc pas de corriger une application mais de construire le socle qui la contraint,
la surveille et la rend reproductible ; le chapitre suivant examine l'état de l'art et justifie les
choix technologiques.
